
\examquestion{2024}{4}{Two-port cavity Langevin equation and photonic qubit encoding}{5}
Given the cavity damping rate \(\gamma=\gamma_1+\gamma_2\),
\[
\dot a=-\frac{\gamma}{2}a+
\sqrt{\gamma_1}A_{{\rm in},1}
+\sqrt{\gamma_2}A_{{\rm in},2}.
\]
\subsection*{(a) Fourier-domain intracavity field}
Take
\(\tilde a(\omega)=\int e^{-i\omega t}a(t)\diff t\),
so \(\widetilde{\dot a}=i\omega\tilde a\).
Collecting terms,
\[
\left(\frac{\gamma}{2}+i\omega\right)\tilde a
=\sqrt{\gamma_1}\tilde A_{{\rm in},1}
+\sqrt{\gamma_2}\tilde A_{{\rm in},2},
\]
\[
\boxed{\tilde a=
\frac{\sqrt{\gamma_1}\tilde A_{{\rm in},1}
+\sqrt{\gamma_2}\tilde A_{{\rm in},2}}
{\gamma/2+i\omega}}.
\]

\subsection*{(b) Field reflected from the first port}
Use the specified boundary condition
\(\tilde A_{{\rm out},1}=
\sqrt{\gamma_1}\tilde a-\tilde A_{{\rm in},1}\).
After substituting,
\[
\boxed{\tilde A_{{\rm out},1}=
\frac{\gamma_1-\gamma/2-i\omega}{\gamma/2+i\omega}
\tilde A_{{\rm in},1}
+\frac{\sqrt{\gamma_1\gamma_2}}{\gamma/2+i\omega}
\tilde A_{{\rm in},2}.}
\]
The second input is the internal loss reservoir and is
required in the \emph{operator} output, even if its
mean coherent amplitude is zero.

\subsection*{(c) Critical coupling at zero Fourier frequency}
For \(\gamma_1=\gamma_2\), we have
\(\gamma_1=\gamma/2\). Therefore the \emph{coherent
reflection coefficient} at \(\omega=0\) is
\[
\boxed{r_1(0)=\frac{\gamma_1-\gamma/2}{\gamma/2}=0}.
\]
A coherent field incident at port 1 then has zero
mean resonant reflected amplitude and, ideally, no
reflected coherent intensity. Energy is lost through
port 2 (the internal-loss channel); unavoidable reservoir
vacuum fluctuations remain in the output operator.

\subsection*{(d) Two dual-rail encodings}
A dual-rail qubit places one photon in exactly one of
two distinguishable optical modes.
\begin{itemize}
\item \textbf{Spatial path}: \(\ket{0_L}=\ket{1_a0_b}\),
\(\ket{1_L}=\ket{0_a1_b}\). Natural for chip-based
interferometers with spatially separate waveguides.
\item \textbf{Polarisation modes}: \(\ket{0_L}=\ket{H}\),
\(\ket{1_L}=\ket V\), i.e.\ occupancy of the \(H/V\)
polarisation modes. Efficient in a single fibre/path
where maintaining two spatial arms is inconvenient.
\end{itemize}
Path encoding is often robust to uniform polarisation
rotations but sensitive to relative path phase;
polarisation encoding avoids distinct paths but can
suffer birefringence and polarisation drift.

\subsection*{(e) Implement a Hadamard physically}
For the logical basis, the Hadamard is
\[
H=\frac1{\sqrt2}\begin{pmatrix}1&1\\1&-1\end{pmatrix}.
\]
In \textbf{path encoding}, a balanced beam splitter
with the appropriate phase convention implements \(H\)
(up to phase shifters). In \textbf{polarisation encoding},
a half-wave plate with its fast axis at \(22.5^\circ\)
to horizontal implements
\(\ket H\mapsto(\ket H+\ket V)/\sqrt2\),
\(\ket V\mapsto(\ket H-\ket V)/\sqrt2\),
up to overall/relative phases set by conventions.
\solcheck{At equal rates, the two-port cavity must
still preserve the combined input/output commutation
relations; discarding port 2 at the operator level
would violate quantum noise balance.}
